\documentclass{article}
\usepackage[UTF8]{ctex}
\usepackage{amsmath}
\usepackage{amssymb}
\usepackage{graphicx}
\usepackage{geometry}
\geometry{a4paper, margin=2cm}

\newcommand{\problemname}[1]{\begin{center}\Large\bfseries #1\end{center}}
\usepackage{enumitem}

\begin{document}
\problemname{幸运的数字}

蒜头君认为有些数字可以给他带来好运气，他一直把 \textbf{首位和末位相乘小于首位和末位相加} 的三位数认为是自己的幸运数。

\subsection*{输入格式}
输入一个数字 $x(100\lex\le999)$

\subsection*{输出格式}
输出为一行，如果是蒜头君的幸运数，则输出 \texttt{YES}，否则输出 \texttt{NO}。

\subsection*{样例输入}
\begin{verbatim}
888
\end{verbatim}

\subsection*{样例输出}
\begin{verbatim}
NO
\end{verbatim}

\begin{itemize}
  \item 首位和末位相乘：$8\times8=64$

  \item 首位和末位相加：$8+8=16$

\end{itemize}
$64>16$，因此 $888$ 不是幸运数。

\subsection*{样例 2 输入}
\begin{verbatim}
124
\end{verbatim}

\subsection*{样例输出}
\begin{verbatim}
YES
\end{verbatim}

\begin{itemize}
  \item 首位和末位相乘：$1\times4=4$

  \item 首位和末位相加：$1+4=5$

\end{itemize}
$4<5$，因此 $124$ 是幸运数。


\end{document}
